% =====================================================================
%  CAPÍTULO 2 · DESARROLLO DEL TEMA
% =====================================================================
\chapter{Desarrollo del tema}
\label{cap:desarrollo}

En este capítulo se desarrolla el análisis central de la monografía, aplicando
el marco teórico del Capítulo~\ref{cap:marco}.

\section{Descripción del problema}
\label{sec:problema}

Presente los hechos, datos o casos que se analizarán, indicando siempre su
procedencia.

\section{Análisis}
\label{sec:analisis}

Desarrolle los argumentos de manera ordenada. Si necesita una expresión
matemática, puede incluirla numerada:

\begin{equation}
  \bar{x} = \frac{1}{n}\sum_{i=1}^{n} x_i
  \label{eq:media}
\end{equation}

donde $\bar{x}$ es la media aritmética y $n$ el número de observaciones
(ecuación~\ref{eq:media}).

\section{Ilustraciones}
\label{sec:ilustraciones}

Para insertar una figura, guarde la imagen en la carpeta \texttt{figuras/} y
use el siguiente bloque (quite el \texttt{\%} para activarlo):

\begin{verbatim}
\begin{figure}[htbp]
  \centering
  \includegraphics[width=0.7\textwidth]{nombre-de-la-imagen.png}
  \caption{Descripción de la figura.}
  \label{fig:ejemplo}
\end{figure}
\end{verbatim}

Recuerde indicar la fuente debajo de cada figura o tabla tomada de otros
autores, siguiendo la norma de citación de su institución.

\section{Discusión parcial}
\label{sec:discusion-parcial}

Interprete lo analizado y conéctelo con los antecedentes expuestos en la
Sección~\ref{sec:antecedentes}.
